\chapter[Introduction]{Introduction: reinforcement learning and the Pacman
  example}
\label{ch:pacman}

Reinforcement learning studies how an \emph{agent} can learn to act in an
\emph{environment} so as to collect as much \emph{reward} as possible, using
only the experience it gathers by acting. Nobody tells the agent which
action is correct: it only observes the consequences of its actions, which
may come much later. This chapter introduces the vocabulary of RL, gives a
map of its methods, and states precisely the rules and rewards of the Pacman
game of the repository, which the following chapters use as a running
example whenever an algorithm applies to it.

\section{The agent--environment loop}

At each time step $t = 0, 1, 2, \dots$ the agent observes the state $S_t$ of
the environment, chooses an action $A_t$, and the environment answers with
a reward $R_{t+1}$ and a new state $S_{t+1}$
(\cref{fig:loop}). The interaction produces a \emph{trajectory}
\begin{equation}
  S_0,\, A_0,\, R_1,\, S_1,\, A_1,\, R_2,\, S_2,\, A_2,\, R_3,\, \dots
\end{equation}
which stops when a \emph{terminal} state is reached; the trajectory is then
called an \emph{episode}. For Pacman an episode is a game.

\begin{figure}[ht]
  \centering
  \begin{tikzpicture}[node distance=2.2cm and 3cm, >=Stealth,
      box/.style={draw, rounded corners, minimum width=3cm,
        minimum height=1cm, thick}]
    \node[box, fill=defblue!10] (agent) {Agent (policy $\pi$)};
    \node[box, fill=pacgreen!10, below=1.6cm of agent] (env)
      {Environment (the game)};
    \draw[->, thick] (agent.east) -- ++(1.6,0) |- node[pos=0.25, right]
      {action $A_t$} (env.east);
    \draw[->, thick] (env.west) -- ++(-1.6,0) |- node[pos=0.25, left]
      {state $S_{t+1}$, reward $R_{t+1}$} (agent.west);
  \end{tikzpicture}
  \caption{The agent--environment loop.}
  \label{fig:loop}
\end{figure}

The behaviour of the agent is described by a \emph{policy} $\pi$, which gives
for each state a probability distribution over the actions; the goal of RL
is to find a policy that maximises the expected sum of the rewards
(\cref{ch:mdp} makes this precise).

\section{Where reinforcement learning is used}

The framework covers any problem of sequential decisions under
uncertainty in which the quality of a behaviour can be measured by
rewards. Some landmarks:
\begin{itemize}[nosep]
  \item \textbf{Games}: backgammon (TD-Gammon, 1992), Atari video games from
    pixels (DQN, 2015), Go, chess and shogi (AlphaGo and AlphaZero,
    2016--2018), StarCraft and Dota~2 (2019).
  \item \textbf{Robotics and control}: locomotion, manipulation, the control
    of plasma in a tokamak, the cooling of data centres; actions are then
    continuous (\cref{ch:continuous}).
  \item \textbf{Recommendation and operations}: choosing what to show to a
    user, pricing, routing; often formulated as bandits (\cref{ch:bandits}).
  \item \textbf{Language models}: the alignment of large language models
    with human preferences (\cref{ch:rlhf}).
\end{itemize}

\section{A map of the methods}
\label{sec:map}

RL algorithms can be classified along a few independent axes, which
organise this course.
\begin{itemize}
  \item \textbf{What is learnt.} \emph{Value-based} methods learn the value
    of states or actions and act greedily with respect to it (Q-learning,
    DQN, \cref{ch:td,ch:dqn}). \emph{Policy-based} methods learn the policy
    directly (policy gradients, \cref{ch:pg}). \emph{Actor-critic} methods
    learn both (A2C, PPO, SAC, \cref{ch:pg,ch:ppo,ch:continuous}).
    \emph{Model-based} methods learn or use a model of the environment to
    plan (\cref{ch:dp,ch:model-based}).
  \item \textbf{Which data.} \emph{On-policy} methods learn from the data of
    the current policy (SARSA, A2C, PPO); \emph{off-policy} methods can learn
    from data generated by other policies (Q-learning, DQN, SAC);
    \emph{offline} methods learn from a fixed dataset without interacting
    with the environment (\cref{ch:imitation}).
  \item \textbf{How targets are formed.} \emph{Monte Carlo} methods use
    complete returns; \emph{temporal-difference} methods \emph{bootstrap}
    from their own estimates (\cref{ch:td}).
  \item \textbf{Which actions.} Discrete actions (Pacman's five moves) or
    continuous ones (the torques of a robot, \cref{ch:continuous}).
  \item \textbf{How many agents.} One agent in a fixed environment, or
    several agents learning together (\cref{ch:multiagent}).
\end{itemize}
Two further themes cut across all of them: \emph{exploration}, the need to
try actions whose value is uncertain (\cref{ch:bandits,ch:exploration}),
and \emph{generalisation}, the need to represent values and policies with
parametrised functions when there are too many states for a table
(\cref{ch:vfa,ch:dqn}).

\section{The running example: a Pacman game}
\label{sec:game}

The game of the repository is a simplified Pacman written from scratch to
test RL algorithms. We describe the original version of the game (called
\code{v1} in the code); \cref{sec:v2} lists the changes of version \code{v2}.

\begin{figure}[ht]
  \centering
  \begin{tikzpicture}[scale=0.55, ghost/.style={rectangle, rounded corners=2pt,
    minimum width=0.32cm, minimum height=0.32cm, inner sep=0pt, draw=black}]
  \fill[black!80] (0,0) rectangle (15,1) (0,-1) rectangle (1,0)
     (7,-1) rectangle (8,0) (14,-1) rectangle (15,0) (0,-2) rectangle (1,-1)
     (2,-2) rectangle (3,-1) (4,-2) rectangle (6,-1) (7,-2) rectangle (8,-1)
     (9,-2) rectangle (11,-1) (12,-2) rectangle (13,-1)
     (14,-2) rectangle (15,-1) (0,-3) rectangle (1,-2)
     (2,-3) rectangle (3,-2) (12,-3) rectangle (13,-2)
     (14,-3) rectangle (15,-2) (0,-4) rectangle (1,-3)
     (2,-4) rectangle (4,-3) (5,-4) rectangle (10,-3)
     (11,-4) rectangle (13,-3) (14,-4) rectangle (15,-3)
     (0,-5) rectangle (1,-4) (14,-5) rectangle (15,-4)
     (0,-6) rectangle (3,-5) (4,-6) rectangle (6,-5) (7,-6) rectangle (8,-5)
     (9,-6) rectangle (11,-5) (12,-6) rectangle (15,-5)
     (0,-7) rectangle (1,-6) (4,-7) rectangle (5,-6)
     (10,-7) rectangle (11,-6) (14,-7) rectangle (15,-6)
     (0,-8) rectangle (1,-7) (2,-8) rectangle (3,-7) (4,-8) rectangle (5,-7)
     (6,-8) rectangle (9,-7) (10,-8) rectangle (11,-7)
     (12,-8) rectangle (13,-7) (14,-8) rectangle (15,-7)
     (0,-9) rectangle (1,-8) (6,-9) rectangle (9,-8)
     (14,-9) rectangle (15,-8) (0,-10) rectangle (15,-9);
  \draw[black!25] (0,-10) grid (15,1);
  \foreach \p in {(1.5,-0.5),(3.5,-0.5),(4.5,-0.5),(5.5,-0.5),(6.5,-0.5),
    (8.5,-0.5),(9.5,-0.5),(10.5,-0.5),(11.5,-0.5),(12.5,-0.5),(13.5,-0.5),
    (1.5,-1.5),(3.5,-1.5),(6.5,-1.5),(8.5,-1.5),(11.5,-1.5),(13.5,-1.5),
    (3.5,-2.5),(4.5,-2.5),(5.5,-2.5),(6.5,-2.5),(7.5,-2.5),(8.5,-2.5),
    (9.5,-2.5),(10.5,-2.5),(11.5,-2.5),(13.5,-2.5),(1.5,-3.5),(4.5,-3.5),
    (10.5,-3.5),(13.5,-3.5),(1.5,-4.5),(2.5,-4.5),(3.5,-4.5),(4.5,-4.5),
    (5.5,-4.5),(6.5,-4.5),(8.5,-4.5),(9.5,-4.5),(10.5,-4.5),(11.5,-4.5),
    (12.5,-4.5),(13.5,-4.5),(3.5,-5.5),(6.5,-5.5),(8.5,-5.5),(11.5,-5.5),
    (1.5,-6.5),(2.5,-6.5),(3.5,-6.5),(5.5,-6.5),(6.5,-6.5),(7.5,-6.5),
    (8.5,-6.5),(9.5,-6.5),(11.5,-6.5),(12.5,-6.5),(1.5,-7.5),(3.5,-7.5),
    (5.5,-7.5),(9.5,-7.5),(11.5,-7.5),(13.5,-7.5),(1.5,-8.5),(2.5,-8.5),
    (3.5,-8.5),(4.5,-8.5),(5.5,-8.5),(9.5,-8.5),(10.5,-8.5),(11.5,-8.5),
    (13.5,-8.5)}
    {\fill[propamber!70!yellow] \p circle (0.12);}
  \fill[pacgreen!80!lime] (1.5,-2.5) circle (0.22);
  \fill[pacgreen!80!lime] (13.5,-6.5) circle (0.22);
  \node[ghost, fill=red!80] at (2.5,-0.5) {};
  \node[ghost, fill=magenta!70] at (12.5,-8.5) {};
  \fill[blue!80] (7.5,-4.5) -- ++(25:0.32) arc (25:335:0.32) -- cycle;
\end{tikzpicture}

  \caption{A position of game \code{v1} on the maze \code{pacman\_map3}:
    Pacman (blue), the two ghosts, the two power cookies (green) and the
    cookies (yellow). The walls are black.}
  \label{fig:board}
\end{figure}

\paragraph{The board.} The maze \code{pacman\_map3} is a grid of
$11 \times 15$ cells whose borders are walls; $77$ cells are open
(\cref{fig:board}). At the start of a game every open cell holds a cookie,
except the cell of Pacman. Two of the cookies are replaced by \emph{power
cookies}. Pacman and the two ghosts start on random cells, with the ghosts
at a distance at least $2$ from Pacman.

\paragraph{The actions.} Pacman has five actions: \code{E}, \code{N},
\code{W}, \code{S} (move by one cell to the east, north, west or south) and
\code{I} (stay idle). Moving into a wall leaves Pacman where he is.

\paragraph{One step of the game.} A call to \code{step(action)} plays one
move of everybody, in this order:
\begin{enumerate}[nosep]
  \item ghost $0$ moves;
  \item Pacman moves; his power credit (see below) decreases by one;
    he eats the cookie of his new cell, if any;
  \item the other ghosts move.
\end{enumerate}
After each of these moves, if Pacman and a ghost are on the same cell, the
ghost eats Pacman (who loses the game) if Pacman has no power, and Pacman
eats the ghost otherwise (the ghost leaves the game).

\paragraph{The power.} Eating a power cookie sets Pacman's \emph{power
credit} to $C = 12$. The credit decreases by one at each move of Pacman, and
Pacman can eat ghosts while it is positive. Since the credit is decreased
\emph{before} the collision test, Pacman with a credit of $1$ who moves onto
a ghost has no power any more at the moment of the collision, and is eaten:
we will meet this ``off-by-one'' situation again in \cref{ch:search}.

\paragraph{The end of a game.} Pacman wins when he has eaten all the
cookies (normal and power ones), and loses when a ghost eats him. Games
longer than a fixed number of moves are \emph{aborted} (1000 moves during
training): this limit is not part of the game itself, and
\cref{ch:mdp} explains why it must be treated differently from a real end.

\paragraph{The ghosts.} The ghosts are the random part of the environment.
Ghost $i$ has an \emph{aggressiveness} $a_i$ ($a_0 = 0.74$, $a_1 = 0.64$).
At each of its moves, when Pacman has no power:
\begin{itemize}[nosep]
  \item with probability $a_i$ the ghost \emph{chases} Pacman: it moves to
    the next cell of a shortest path to Pacman (computed with the A$^\star$
    algorithm);
  \item otherwise it moves to a cell chosen uniformly among the cells it can
    reach in one move, including its own cell.
\end{itemize}
When Pacman has power, the ghosts \emph{flee}: a ghost moves to a reachable
cell chosen with a probability proportional to its Euclidean distance to
Pacman, and never to Pacman's cell.

\begin{remark}{Random moves of the ghosts}{ghost-heading}
  The code of the ghosts intends to give a weight $1.5$ to the move that
  keeps the current direction of a ghost, but the comparison it uses (a
  position with a direction letter) is never true: the random moves are
  uniform, as described above. A measurement over many games confirms it
  (a ghost keeps its direction $31.7\%$ of the time, against $31.2\%$ for a
  uniform choice). As a consequence, the ghosts' moves depend only on the
  current positions and on Pacman's power, which matters in
  \cref{ch:mdp}. Game \code{v2} (\cref{sec:v2}) fixes the comparison: its
  ghosts keep their direction $39.6\%$ of the time in the same measurement,
  as predicted by the weight $1.5$, while the games of \code{v1} are
  unchanged.
\end{remark}

\section{The rewards}
\label{sec:rewards}

The reward $R_{t+1}$ of a step is the sum of the rewards of the events of
the step, listed in \cref{tab:rewards}. Its design reflects what the
author of the game wanted Pacman to learn: eat the cookies, use the power
to eat the ghosts, and avoid dying.

\begin{table}[ht]
  \centering
  \begin{tabular}{@{}lll@{}}
    \toprule
    Event & Reward & Condition \\
    \midrule
    Living cost (each move of Pacman) & $-0.15$ &
      Pacman has no power, or no ghost is left \\
    Staying in place & $-0.75$ &
      instead of the living cost, if Pacman has \\
    & & power or no ghost is left \\
    Eating a cookie & $+0.8$ &
      Pacman has no power, or no ghost is left \\
    Eating a power cookie & see \eqref{eq:power-reward} &
      $+0.8$ if no ghost is left \\
    Eating a ghost & $+20$ & \\
    Being eaten & $-35$ & \\
    Winning & $+30 - 5\,g$ & $g$ = number of ghosts left \\
    \bottomrule
  \end{tabular}
  \caption{Rewards of game \code{v1} (\code{pacman\_training.py}).}
  \label{tab:rewards}
\end{table}

The reward for a power cookie depends on the situation. If $c$ is Pacman's
power credit when he eats it (after the decrease of his move) and $d \geq 1$
is the Manhattan distance to the nearest ghost,
\begin{equation}
  \label{eq:power-reward}
  r_{\text{power}} = 1.3 \left( 1 - \frac{c}{C} \right) \frac{C}{d}.
\end{equation}
Eating a power cookie is thus worth more when Pacman has little power left
and when a ghost is close: $r_{\text{power}} = 1.3 \times 12 = 15.6$ with no
power and a ghost at distance $1$, but only $1.3$ with a ghost at distance
$12$.

\begin{pacman}{An optimistic estimate of the score}{score}
  A game won without power and without eating a ghost, in $n$ moves, has a
  total reward of about $0.8 \times 76 - 0.15\,n + 30 - 10$. With
  $n = 100$ this is $60.8 - 15 + 20 = 65.8$. Each ghost eaten adds $20 + 5$
  (its reward and the penalty avoided), which is why the trained agents
  hunt the ghosts: the reward per game of the agent of the repository is
  about $115$ at the end of its training (\cref{ch:ppo}).
\end{pacman}

\section{What the agent observes}
\label{sec:observation}

The neural networks of the agents do not receive the Python objects of the
game but an image of it, the \emph{feature planes} of
\code{PacmanEnvironment.feature\_image}: an array of shape
$(5 + n_g) \times 9 \times 13$ for $n_g$ ghosts (the borders are removed),
with one binary plane for each kind of object:
\begin{enumerate}[nosep, label=\arabic*.]
  \item the walls;
  \item the cookies;
  \item the power cookies;
  \item Pacman;
  \item[5--6.] one plane per ghost (they have different aggressiveness);
  \item[7.] a constant plane equal to Pacman's power $c / C \in [0, 1]$.
\end{enumerate}
The agents stack the last two images (\emph{frame stacking}), so that the
input has $2 \times 7 = 14$ channels. \cref{ch:mdp} discusses why the image
is (almost) a complete description of the state of the game, and what the
second frame adds.

\section{Why reinforcement learning?}

Pacman illustrates the three difficulties that distinguish RL from
supervised learning.
\begin{itemize}
  \item \textbf{No teacher.} There is no dataset of correct moves: the
    agent only receives rewards, and must discover by itself that, for
    example, waiting near a power cookie for a ghost to come close is a
    good idea.
  \item \textbf{Delayed consequences.} A move into a dead end can be fatal
    five moves later; the reward $-35$ arrives long after the mistake. The
    agent must assign the credit (or the blame) of a reward to the actions
    that caused it.
  \item \textbf{Exploration.} The data depend on the agent's own
    behaviour: an agent that never eats a power cookie never learns that
    eating ghosts pays. The agent must balance the exploitation of what it
    knows with the exploration of what it does not.
\end{itemize}

\section{Game v2}
\label{sec:v2}

On game \code{v1} the trained agents win more than $96\%$ of the games and
usually eat both ghosts early, after which the game is only about cookies.
Game \code{v2} makes the whole game matter: a bigger maze
(\code{pacman\_map4}, $17 \times 19$ cells, $163$ open), three ghosts, four
power cookies, a power of $10$ moves, a reward of $5$ for eating a ghost,
ghosts that \emph{come back} (an eaten ghost reappears $15$ moves after
being eaten, and not before Pacman's power has ended, on a random cell at
least $6$ moves away from Pacman), and ghosts whose random moves prefer
their current direction (\cref{rem:ghost-heading}). Unless stated
otherwise, the examples of this course use game \code{v1}, on which all the
results were obtained.

\section{Plan of the course}

\begin{itemize}[nosep]
  \item \textbf{Foundations.} \cref{ch:bandits} studies exploration in the
    simplest setting, bandits. \cref{ch:mdp} gives the framework (Markov
    decision processes, Bellman equations), \cref{ch:dp} solves it when the
    model is known, \cref{ch:td} learns from experience instead.
  \item \textbf{Function approximation.} \cref{ch:vfa} studies linear value
    functions and the stability of TD learning; \cref{ch:dqn} moves to deep
    networks (DQN and its extensions).
  \item \textbf{Policy optimisation.} \cref{ch:pg} derives policy gradients
    and actor-critic methods, \cref{ch:ppo} trust regions (TRPO, PPO),
    \cref{ch:continuous} continuous actions (DDPG, TD3, SAC).
  \item \textbf{Planning and data.} \cref{ch:model-based} covers model-based
    RL and search (Dyna, MCTS, AlphaZero, MuZero), \cref{ch:imitation}
    learning from demonstrations and from fixed datasets.
  \item \textbf{Further topics.} \cref{ch:exploration} returns to
    exploration in large environments, \cref{ch:multiagent} to several
    agents, \cref{ch:rlhf} to learning from human preferences.
\end{itemize}
The Pacman examples use the code of the repository; code excerpts show
how the equations are implemented. The chapters on continuous actions,
multi-agent RL and human feedback use other examples, as these settings do
not arise in the game.
